% --- Άσκηση 13114 (13114.tex) ---
\Askhsh{13114}{4}
\begin{ylh}{2}
    \item 2.3. Απόλυτη Τιμή Πραγματικού Αριθμού
    \item 3.1. Εξισώσεις 1ου Βαθμού
    \item 4.1. Ανισώσεις 1ου Βαθμού
\end{ylh}
Δίνεται η συνάρτηση $f(x) = \dfrac{|x - 1| - |3 - 3x| + |2x - 4|}{2}$.
\begin{alist}
\item Να δείξτε ότι $f(x) = d(x,2) - d(x,1)$.
\monades{9}
\item Αν τα σημεία $A$ και $B$ παριστάνουν στον άξονα των πραγματικών αριθμών τους αριθμούς 1 και 2, να διατυπώσετε γεωμετρικά το ζητούμενο της εξίσωσης $f(x) = 0$ και να προσδιορίσετε τη λύση της.

\monades{8}
\item Να λύσετε την εξίσωση $f(x) = 0$.
\monades{8}
\end{alist}
